% Generated draft: requires mathematical and source review.
\documentclass[11pt]{article}
\usepackage[margin=25mm]{geometry}
\usepackage{fontspec}
\setmainfont{Latin Modern Roman}
\usepackage{amsmath,amssymb,mathtools,unicode-math}
\setmathfont{Latin Modern Math}
\usepackage{enumitem,needspace,xurl,hyperref}
\setlength{\parindent}{0pt}
\setlength{\parskip}{5pt}
\setlength{\emergencystretch}{3em}
\newcommand{\N}{\mathbb{N}}
\newcommand{\Z}{\mathbb{Z}}
\newcommand{\Q}{\mathbb{Q}}
\newcommand{\R}{\mathbb{R}}
\newcommand{\C}{\mathbb{C}}
\newcommand{\F}{\mathbb{F}}
\newcommand{\A}{\mathbb{A}}
\newcommand{\PP}{\mathbb{P}}
\newcommand{\E}{\mathbb{E}}
\newcommand{\eps}{\varepsilon}
\newcommand{\dd}{\,\mathrm{d}}
\begin{document}
\section*{UnsolvedMath ID: 30003630}
\textbf{Source problem code:} catalogue entry
\subsection{Definitions and mathematical statement}
For the conventional interpretation of the terminology, a Kodaira fibration is a surjective holomorphic submersion \(f\colon S\to C\) from a smooth compact connected complex surface to a smooth compact connected complex curve, with connected fibres, such that \(g(C)\geq 2\), every fibre has genus at least \(2\), and the family is non-isotrivial, meaning that its fibres are not all mutually biholomorphic. The supplied record does not explicitly state these hypotheses; this definition makes the conventional interpretation explicit and requires confirmation against the source.

Under this interpretation, let \(B\) and \(F_b\) be smooth compact connected complex curves of genus at least \(2\). Consider a short exact sequence of groups
\[
\mathcal E\colon\quad 1\longrightarrow\pi_1(F_b)\xrightarrow{\iota}G\xrightarrow{p}\pi_1(B)\longrightarrow1,
\]
where \(\iota\) is injective, \(p\) is surjective, and \(\operatorname{im}(\iota)=\ker(p)\). Fundamental groups are taken with compatible basepoints, suppressed from the notation.

\textbf{Problem, following the corrected registry wording.} Characterize the extensions \(\mathcal E\) for which there exist a Kodaira fibration \(f\colon S\to C\) and an abstract group isomorphism
\[
G\cong\pi_1(S).
\]

\textbf{Unresolved scope of realization.} The displayed condition concerns the abstract middle group. Realization of the given extension is a more specific requirement: there must also be a point \(c\in C\) and isomorphisms making the following diagram commute, with the bottom row induced by \(f\):
\[
\begin{array}{ccccccccc}
1&\longrightarrow&\pi_1(F_b)&\xrightarrow{\iota}&G&\xrightarrow{p}&\pi_1(B)&\longrightarrow&1\\
&&\downarrow\scriptstyle\cong&&\downarrow\scriptstyle\cong&&\downarrow\scriptstyle\cong&&\\
1&\longrightarrow&\pi_1(f^{-1}(c))&\longrightarrow&\pi_1(S)&\xrightarrow{f_*}&\pi_1(C)&\longrightarrow&1.
\end{array}
\]
The supplied wording does not determine whether this compatibility is required, or whether \(C\) and \(f^{-1}(c)\) must be biholomorphic to the specified curves \(B\) and \(F_b\). These requirements must not be silently identified with abstract group realization.

The original wording asks for the fundamental group of a ``Kodaira surface,'' whereas the corrected wording asks for that of a Kodaira fibration. Here the corrected wording supplies the target; the original terminology remains unresolved because ``Kodaira surface'' also has a distinct standard meaning in complex-surface classification.
\subsection{Short English statement}
Determine which extensions of a curve fundamental group by another curve fundamental group have a middle group isomorphic to the fundamental group of a Kodaira fibration. The source leaves unresolved whether the specified extension itself, and the specified complex curves, must also be realized.
\subsection{Sources}
\begin{itemize}
\item Supplied local UnsolvedMath registry record, permanent ID \(30003630\), exact source problem code \texttt{catalogue entry}, title \emph{Group Extensions Realized by Kodaira Fibrations}, in the collection \emph{Oberwolfach Reports Open Problems}. The proposed year is recorded as \(2017\); no individual proposer is supplied.
\item Supplied report citation: \emph{Mini-Workshop: Interactions between Low-dimensional Topology and Complex Algebraic Geometry (2018)}, Oberwolfach Reports, pp.\ 2987--3033, DOI \texttt{10.4171/owr/2017/50}. \url{https://doi.org/10.4171/owr/2017/50}. The literature-source label is ``Oberwolfach Report 50/2017.'' The supplied citation's year \(2018\) differs from the report label and proposed year \(2017\); this discrepancy is unresolved.
\item Supplied literature source labelled ``Vidussi'': \url{https://arxiv.org/abs/1706.03197}. No paper title or further bibliographic details are supplied.
\item Supplied literature assessment, checked on \texttt{2026-08-21}: the registry labels the problem \texttt{partially\_solved} and states that Vidussi derives necessary restrictions and gives extensions that cannot be realized by Kodaira fibrations, while no complete characterization of realizable extensions is available. This is supplied dated curation derived from the OWR record, not an independently verified assessment. The record expressly requests independent verification of time-sensitive status and cited claims.
\end{itemize}
\end{document}
